\section{INTRODUCTION}

Automated unit test generation plays a critical role in improving software reliability and reducing developer effort. Traditional techniques, such as symbolic execution and search-based software testing (SBST), have made substantial progress in systematically exploring input spaces and maximizing structural coverage~\cite{cadar2008klee, evosuite, mcminn2011}. However, these approaches often face fundamental limitations in scalability and in handling complex program semantics. In particular, they struggle to capture implicit preconditions that arise from intricate object states and inter-procedural interactions, where the conditions required to reach a target path are established outside the local method scope through sequences of prior method calls~\cite{baldoni2018surveysymbolicexecutiontechniques}.

Recent advances in large language models (LLMs) have introduced a new paradigm for unit test generation. By leveraging their strong code understanding and reasoning capabilities, LLM-based approaches can directly synthesize executable test cases from source code, significantly lowering the barrier to automated testing~\cite{chen2024, ryan2024, wang2024c, gu2025a, liao2025}. Moreover, iterative workflows that incorporate compilation and execution feedback have further improved the validity and coverage of generated tests~\cite{pizzorno2024}.

Despite these promising results, existing LLM-based methods still exhibit a fundamental limitation: they lack precise reasoning about \emph{implicit program constraints}, particularly those that span across multiple methods. Most approaches rely on prompting the LLM with the entire method body, optionally augmented with control-flow or type information. However, in real-world object-oriented programs, the conditions required to execute a target branch are often not locally visible within the method under test. Instead, they depend on object states established through prior method calls or initialization sequences, which are external to the focal method. As a result, LLMs frequently generate tests that are syntactically correct but fail to satisfy the hidden preconditions necessary for reaching uncovered branches~\cite{ryan2024}.

Furthermore, providing the full method body as context introduces substantial noise. Irrelevant statements consume the limited context window and obscure the critical dependencies that govern branch reachability. This not only increases the cognitive burden on the LLM but also reduces its ability to infer correct execution conditions. Consequently, even when guided by path information, existing approaches often fail to cover branches that require precise state construction, as LLM performance has been shown to degrade significantly as input context length increases, even in the absence of distracting information~\cite{du2025contextlengthhurtsllm, wang2026intelligencedegradationlongcontextllms}.

To address these challenges, we propose \ourmethod{}, a novel LLM-driven framework for automated unit test generation that tightly integrates program analysis with language model reasoning. The key idea is to reduce the difficulty of test generation from two complementary perspectives: \emph{making hidden constraints explicit} and \emph{eliminating irrelevant context}.

First, we introduce \textit{Constraints-Hints}, a mechanism that extracts path conditions from the control-flow graph and translates them into explicit, human-interpretable input requirements. Instead of requiring the LLM to infer complex branch predicates from raw source code, our approach directly exposes the necessary conditions for covering a target path, enabling more accurate and guided test generation.

Second, we propose a \textit{coverage-driven backward slicing} technique that identifies the minimal subset of program statements relevant to an uncovered branch. By focusing only on statements that influence the target condition—including cross-method dependencies such as field initialization and setter methods—the resulting slice provides a concise and semantically meaningful context for the LLM.

These two components are integrated into an iterative framework that combines test generation, execution, and automated repair. Starting from an initial test suite, the framework continuously identifies uncovered paths, derives constraint hints, and generates candidate tests. When coverage improvement stagnates, backward slicing is triggered to provide more precise guidance. This iterative feedback loop enables sustained coverage growth and effective exploration of hard-to-reach program paths.

We evaluate \ourmethod{} on a set of real-world Java projects and compare it against state-of-the-art LLM-based and search-based baselines. The results show that our approach consistently achieves higher branch coverage and is particularly effective in scenarios involving complex object state dependencies. Further analysis demonstrates that both constraint hints and backward slicing are essential for guiding LLMs toward generating semantically correct and coverage-effective tests.

In summary, this paper makes the following contributions:

\begin{itemize}[leftmargin=*]
\item We identify a key limitation of existing LLM-based test generation approaches in handling implicit, cross-method constraints.
\item We propose \textit{Constraints-Hints}, a mechanism that makes path conditions explicit and guides LLM reasoning.
\item We introduce a coverage-driven backward slicing technique to eliminate irrelevant context and expose critical dependencies.
\item We design an iterative framework that integrates generation, execution feedback, and repair for sustained coverage improvement.
\item We empirically demonstrate the effectiveness of our approach on real-world software systems.
\end{itemize}